\section{Failure Taxonomy for AI Security Agents}

Understanding \textit{how} agents fail is a prerequisite for measuring
reliability. We propose a three-axis taxonomy that decomposes each
failure event along three orthogonal dimensions: what caused it
(\textit{origin}), what it looked like (\textit{manifestation}), and
what it cost (\textit{impact}). Unlike prior fault taxonomies developed
for conventional software~\cite{avizienis2004}, this taxonomy is
designed specifically for autonomous AI security agents, where
failures are not bugs to be patched but operational risks that affect
real-world security outcomes. Table~\ref{tab:taxonomy_mag} summarises
the three axes.

\begin{table}[t]
\caption{Three-Axis Failure Taxonomy for AI Security Agents}
\label{tab:taxonomy_mag}
\centering
\footnotesize
\setlength{\tabcolsep}{3pt}
\begin{tabular}{p{0.85cm} p{2.5cm} p{4.0cm}}
\toprule
\textbf{Code} & \textbf{Class} & \textbf{Description} \\
\midrule
\multicolumn{3}{l}{\textit{Axis 1: Origin (root cause)}} \\
O1 & Epistemic Deficiency & Lacks required threat knowledge \\
O2 & Skill Boundary & Task requires an unlearned procedure \\
O3 & Capability Deprivation & Required tool or API unavailable \\
O4 & Reasoning Corruption & Inference produces invalid output \\
O5 & Context Overload & Input exceeds processing capacity \\
O6 & Adversarial Perturbation & Malicious content degrades output \\
O7 & Orchestration Failure & Multi-agent coordination breakdown \\
O8 & Temporal Drift & Knowledge has become stale \\
\midrule
\multicolumn{3}{l}{\textit{Axis 2: Manifestation (observable signal)}} \\
M1 & Hallucination & Asserts verifiably false facts \\
M2 & Omission & Fails to detect a present threat \\
M3 & Misclassification & Wrong threat family or category \\
M4 & Confabulation & Coherent reasoning from false premises \\
M5 & Paralysis & No output or persistent retry loop \\
M6 & Overclaiming & High confidence on weak evidence \\
M7 & Attribution Error & Wrong threat actor or campaign \\
M8 & Cascade Failure & Error propagates across agents \\
\midrule
\multicolumn{3}{l}{\textit{Axis 3: Impact (operational consequence)}} \\
I1 & False Negative\hfill\textbf{*} & Threat undetected; attacker persists \\
I2 & False Positive & Benign activity flagged \\
I3 & Delayed Response & Action issued past SLA \\
I4 & Wrong Remediation\hfill\textbf{*} & Incorrect containment executed \\
I5 & Intelligence Pollution\hfill\textbf{*} & Fabricated IOCs propagate to feeds \\
I6 & Trust Erosion & Analyst confidence reduced \\
I7 & Compliance Violation & Regulatory threshold breached \\
\bottomrule
\multicolumn{3}{l}{\footnotesize \textbf{*}\;Critical impact category.} \\
\end{tabular}
\end{table}

Two distinctions matter for practitioners. First, \textit{Hallucination}
(M1) and \textit{Confabulation} (M4) are not synonyms. Hallucination
asserts atomic facts that are verifiably false against a closed-world
corpus; confabulation produces internally coherent reasoning founded
on non-existent premises. Confabulation is more dangerous because it
is harder to detect by automated checks and more persuasive to
analysts reviewing the agent's reasoning trace. Second,
\textit{Intelligence Pollution} (I5) is qualitatively unlike the other
impacts: a single hallucinated IOC pushed to a STIX/TAXII feed
propagates beyond the originating organisation to the entire
information-sharing community.

The power of the three-axis decomposition is its diagnostic
composability. A practitioner who observes a cluster of M1 events can
trace back to O1 (knowledge deficiency) or O5 (context overload) as
the most probable origins, and forward to I1 (false negative) or I5
(intelligence pollution) as the likely consequences. Each origin maps
to a distinct operational action: O1 calls for knowledge refresh, O5
for input preprocessing, O3 for tool provisioning.

\textbf{Worked example: Intelligence Pollution.}
A pipeline correctly identifies a phishing campaign, but the enrichment
stage invents three IOC URLs by extrapolating a domain pattern; a
downstream automation pushes them to the shared feed, and within minutes
partner SOCs block legitimate subdomains of a real service. The trace
decomposes cleanly: origin O1 (missing TTP chain), manifestation M1
(hallucinated facts), impact I5 (blast radius beyond the originating
organisation), and it motivates per-IOC confidence gating,
confidence-proportional feed propagation, and human review for feed
writes, each traceable to a distinct axis.
